\section{Proofs and numerical protocol}
\label{app:vsprotocol}

\subsection{The backward equation for the expected remaining variance}

\Cref{sec:vsthree} stated that
$\mathcal{W}(s,y)=\E[\int_s^{\tau}v(r,Y_r)\,\dd r\,|\,Y_s=y]$
solves the backward problem \eqref{eq:vssource}.  Here is the
derivation, in two steps the chapter has already rehearsed.

First, a martingale.  Consider the running quantity
\[
M_s=\int_0^{s}v(r,Y_r)\,\dd r+\mathcal{W}(s,Y_s),
\]
the variance accumulated so far plus the expectation of what remains.
Write $\mathcal{F}_s$ for the information available at time $s$.
Conditioning the total $\int_0^{\tau}v\,\dd r$ on $\mathcal{F}_s$
splits it into exactly these two pieces --- the first is known at
time $s$, and the second is its conditional expectation by the Markov
property of the diffusion.  So
$M_s=\E[\int_0^{\tau}v\,\dd r\,|\,\mathcal{F}_s]$: every $M_s$ is a
conditional expectation of one fixed random variable, and such a
family is a martingale by the tower property --- conditioning $M_t$
on the smaller information set $\mathcal{F}_s$ ($s<t$) collapses it
back to $M_s$, which is exactly the martingale identity
$\E[M_t\,|\,\mathcal{F}_s]=M_s$.

Second, It\^o.  Assume $\mathcal{W}$ is smooth and apply It\^o's
formula to $\mathcal{W}(s,Y_s)$ with $\dd Y=\sqrt{v}\,Y\dd W$, using
$(\dd W)^2=\dd s$ as in \cref{sec:vsnumber}:
\[
\dd\mathcal{W}(s,Y_s)
=\Big[\partial_s\mathcal{W}
+\tfrac12\,v(s,Y_s)\,Y_s^{2}\,\partial_{yy}\mathcal{W}\Big]\dd s
+\partial_y\mathcal{W}\,\sqrt{v}\,Y_s\,\dd W_s ,
\]
so that
\[
\dd M_s
=\Big[v+\partial_s\mathcal{W}
+\tfrac12\,v\,Y_s^{2}\,\partial_{yy}\mathcal{W}\Big]\dd s
+\partial_y\mathcal{W}\,\sqrt{v}\,Y_s\,\dd W_s .
\]
A martingale has no drift, so the bracket must vanish at every state
the process can reach --- which is the PDE of \eqref{eq:vssource}.
The terminal condition $\mathcal{W}(\tau,\cdot)=0$ holds because
nothing remains at expiry, and evaluating the definition at
$(s,y)=(0,1)$ gives
$\mathcal{W}(0,1)=\E[\int_0^{\tau}v\,\dd r]=w_{\rm vs}$.

\subsection{Numerical protocol}

Every empirical number in the chapter comes from one deterministic
generator reading the book's frozen snapshot; none is typed by hand.
The conventions, in the order the chapter uses them.

\emph{Fits.}  The running node is SPY December 2026.  Its three
family fits are Chapter~3's comparison protocol applied verbatim
(\cref{sec:comparison}): the same mid quotes, equal weights, each
family's reference settings.  The gallery and term-structure figures
(\cref{fig:vsshare}b and~\ref{fig:vsterm}) instead read the
snapshot's \emph{stored} surface fits --- the calibration the surface
was published with, under the fit recipe of \cref{sec:calibration}
--- because those are calendar-coupled across expiries, which is
what a term-structure exhibit should show.

\emph{The strike-side quadrature.}  Equation \eqref{eq:vsrepl} is
evaluated by a trapezoid on $k\in[-6,6]$ with $4001$ points --- a
finer copy of the $801$-point grid the reference implementation uses
inside its optimizer loop; the total variance inside the Black
call is floored at $10^{-12}$.  At these settings the quadrature
error is far below every number quoted.  The law-side integral rides
each slice's own fixed log-odds grid.

\emph{The field-side audit.}  The sheet is Chapter~4's whole-surface
SPY calibration, unchanged.  Both of its var-swap routes run on an
audit lattice with the calibration lattice's strike step, widened to
strike ratio $12$ so the static replication sees the entire wing the
extrapolated field implies; the put leg is cut off below strike
ratio $0.01$ (moving the cutoff to $0.001$ changes nothing at the
quoted precision).  Both routes are first-order in the time step, so
the audit reports the pair of reads at the march's own step and at an
eightfold refinement; the gap halves with the step, as first order
promises.

\emph{The envelope.}  The drawn cone is built from the Black prices
of the last two call-side quote \emph{mids}.  The per-family
clearances are measured against each family's \emph{own} cone ---
the same construction fed that family's fitted call values at the
same two strikes --- because \cref{prop:vscone} constrains a
completion relative to its own belly; nonnegativity of all six
clearances is the proposition confirmed numerically.  Clearances are
measured on $\bar k+0.05\le k\le1.8$; beyond $1.8$ the flattest
wing's call price falls under the double-precision floor and its
Black inversion is noise.

\emph{The hat.}  The counterexample of \cref{fig:winglimit} adds one
kernel of \cref{eq:hatdef} with half-width $0.55$ and steepness $4$
in standardized moneyness, centered $1.4$ standardized units past
the last put-side quote.  Its amplitude is the smallest entry of a
fixed ladder ($5$, $10$, $20$, $35$, $55$ percent of the base
variance at the center) that drives the minimum of $g_{\rm D}$ below
$-0.05$ while total variance stays positive.  Realized wing slopes
are measured by finite differences of $w$ at $|k|=8$.
